\input{../tex-inputs/latex-leseansicht-vorspann}

\section[Gustav Schwarzkopf an Arthur Schnitzler, 28. 11. 1893]{L04277 Gustav Schwarzkopf an Arthur Schnitzler, 28. 11. 1893}
\nopagebreak\mylabel{L04277v}
\rehead{ }\normalsize\beginnumbering\briefempfaengerindex{Schnitzler, Arthur@\textsc{Schnitzler, Arthur}!zzzSchwarzkopf, Gustav@\emph{von Gustav Schwarzkopf}!1893-11-281@{28. 11. 1893}|(be}
\toendnotes[C]{\smallbreak\pagebreak[2]}
\correspDesc{Versand  durch Gustav Schwarzkopf am 28. 11. 1893 in Wien
\newline{}Übermittlung  am 28. 11. 1893 in Wien
\newline{}Zustellung  am 28. 11. 1893 in Wien
\newline{}Erhalt  durch Arthur Schnitzler am 28. 11. 1893 in Wien}\toendnotes[C]{\smallbreak}
\Standort{CUL, Schnitzler, B 96a.}
\physDesc{Postkarte, 308 Zeichen
\newline{}Handschrift: schwarze Tinte, deutsche Kurrent
\newline{}Versand: 1) Stempel: »\nobreak{}\oindex{I., Innere Stadt@\textbf{I., Innere Stadt}, \emph{Verwaltungsgebiet}|pwk}Wien 1/1 1, 28. 11. 93, 10–1V\nobreak{}«.   2) Stempel: »\nobreak{}\oindex{IX., Alsergrund@\textbf{IX., Alsergrund}, \emph{Verwaltungsgebiet}|pwk}Wien 9/3 72, 28. 11. 93, 1.N, bestellt\nobreak{}«. }\toendnotes[C]{\smallbreak}\pstart{}{\pb}Herrn \textsc{Dr
                     Arthur Schnitzler}\pend{}\pstart{}\textsc{Wien\oindex{Wien@\textbf{Wien}, \emph{Verwaltungsgebiet}|pw}}\pend{}\pstart{}IX Frankgaſſe 1\oindex{Wien@\textbf{Wien}!IX., Alsergrund@\textbf{IX., Alsergrund}!Frankgasse 1@\textbf{Frankgasse 1}, \emph{Wohngebäude}|pw}.\pend{}{\bigskip}\vspace{1em}
\pstart
           \raggedleft{}{\pb}28. XI. 93\pend
           
\pstart{}Lieber Freund!\pend\vspace{0.5em}
\pstart
           Ich{ }ſtehe ganz zu Ihrer Verfügung und werde mich morgen Vormittag ins
               Volkstheater\oindex{Wien@\textbf{Wien}!I., Innere Stadt@\textbf{I., Innere Stadt}!Burgtheater@\textbf{Burgtheater}, \emph{Theater}|pw}\eventindex{Volkstheater@\textbf{Volkstheater}!Probe von Das Märchen, 29.11.1893@Probe von Das Märchen, 29.11.1893|pwv} einfinden. Ich fürchte nur, daß Ihnen
               meine »Autorität« in der Sache wenig nützen. wird.\pend
           
\pstart
           Herzlichſt{\\[\baselineskip]}Ihr{\\[\baselineskip]}\spacefill\mbox{GustavSch}\pend
           \leftskip=0em{}
\pstart
           \noindent{}Die \label{K_L04277-1v}\edtext{Karte\eventindex{Raimund-Theater@\textbf{Raimund-Theater}!Premiere von Die gefesselte Phantasie, 28.11.1893@Premiere von Die gefesselte Phantasie, 28.11.1893|pwv} in’s Raimundtheater\orgindex{Raimund-Theater@Raimund-Theater|pw}}{\lemma{\textnormal{\emph{Karte in’s Raimundtheater}}}\Cendnote{\textnormal{Am 28. 11. 1893 eröffnete
                     das \emph{Raimundtheater}\orgindex{Raimund-Theater@Raimund-Theater|pwk}. Gegeben wurde  \emph{Die gefesselte Phantasie}\pwindex{Raimund, Ferdinand 1.\,6.\,1790 Wien – 5.\,9.\,1836 Pottenstein@\textsc{Raimund, Ferdinand} (1.\,6.\,1790 Wien – 5.\,9.\,1836 Pottenstein), \emph{Schauspieler, Dramatiker}!gefesselte Phantasie. Original-Zauberspiel mit Gesang in zwei Aufzügen@\strich\emph{Die gefesselte Phantasie. Original-Zauberspiel mit Gesang in zwei Aufzügen}|pwk} von Ferdinand Raimund\pwindex{Raimund, Ferdinand 1.\,6.\,1790 Wien – 5.\,9.\,1836 Pottenstein@\textsc{Raimund, Ferdinand} (1.\,6.\,1790 Wien – 5.\,9.\,1836 Pottenstein), \emph{Schauspieler, Dramatiker}|pwk}.
                     Schnitzler nahm daran teil.}}}\label{K_L04277-1} haben Sie doch erhalten?\pend
           \selectlanguage{ngerman}\endnumbering\briefempfaengerindex{Schnitzler, Arthur@\textsc{Schnitzler, Arthur}!zzzSchwarzkopf, Gustav@\emph{von Gustav Schwarzkopf}!1893-11-281@{28. 11. 1893}|)be}\mylabel{L04277h}
\begin{anhang}
\end{anhang}\newcommand{\dateiname}{L04277}\newcommand{\titel}{Gustav Schwarzkopf an Arthur Schnitzler, 28. 11. 1893}\newcommand{\editorInnen}{Herausgegeben von Jahnke, SelmaMüller, Martin Anton}\input{../tex-inputs/latex-leseansicht-abspann}
